\documentclass[11pt,a4paper]{article}
\usepackage[utf8]{inputenc}
\usepackage[margin=1in]{geometry}
\usepackage{hyperref}
\usepackage{xcolor}
\usepackage{parskip}

% Remove page numbering
\pagestyle{empty}

% Define colors for links
\hypersetup{
    colorlinks=true,
    linkcolor=blue!70!black,
    urlcolor=blue!70!black,
}

\begin{document}

% Compact header
\noindent
\begin{minipage}[t]{0.5\textwidth}
    \textbf{Samuel Chau}\\
    \href{mailto:schau22@pm.me}{schau22@pm.me}\\
    +61 413 867 624\\
    \href{https://github.com/santiagosayshey}{github.com/santiagosayshey}
\end{minipage}
\hfill
\begin{minipage}[t]{0.45\textwidth}
    \raggedleft
    \today\\[10pt]
    Graduate Program Team\\
    SA Water\\
    Adelaide, South Australia
\end{minipage}

\vspace{20pt}

\noindent Dear SA Water Graduate Program Team,

I'm writing to apply for the 2025 Graduate Program in IT/Cybersecurity. As a final-year Computer Science student at the University of Adelaide, I'm excited about the opportunity to apply my technical skills to infrastructure that actually matters – keeping water flowing to millions of South Australians.

Over the past two years, I've been building and maintaining Dictionarry, an open-source platform that helps people automate their media servers. What started as a personal project to solve my own frustrations has grown to over 300,000 downloads and a community of 1000+ active users. While media servers and water infrastructure might seem worlds apart, the core challenges are surprisingly similar: building reliable systems that people depend on daily, securing critical data, and making complex technology accessible to everyone.

Through this project, I've gained hands-on experience that directly aligns with SA Water's needs:

\textbf{Security-First Thinking}: My cybersecurity major isn't just academic – I've implemented it in practice. Whether it's building end-to-end encryption for my OMesh messaging platform or securing Dictionarry's infrastructure for 20,000 active users, I understand that security isn't an afterthought when people depend on your systems.

\textbf{Data at Scale}: I've built ranking algorithms and manage databases that handle thousands of daily queries. This practical experience in data analytics and system optimization would translate well to SA Water's focus on data-driven decision making and operational efficiency.

\textbf{Real Leadership Experience}: Managing an open-source project has taught me more about leadership than any course could. I coordinate contributors across time zones, mediate technical disagreements, and make architectural decisions that affect thousands of users. I've learned that good leadership means building consensus, not just writing code.

What excites me most about SA Water is the scale of impact. My current work helps people organize their entertainment; working with SA Water would mean contributing to infrastructure that's genuinely essential. The graduate program's rotation structure particularly appeals to me – I've seen firsthand how exposure to different technical challenges accelerates learning and helps identify where you can add the most value.

I'm also drawn to SA Water's commitment to innovation and listening to fresh perspectives. In the open-source world, the best ideas often come from unexpected places. I'd love to bring that collaborative, innovation-focused mindset to help modernize and secure SA Water's digital infrastructure.

Thank you for considering my application. I look forward to discussing how I can contribute to SA Water's mission.

\vspace{20pt}
\noindent
Sincerely,\\
Samuel Chau

\end{document}